\noindent\textbf{Augmentation policies for object detection.} Automatic augmentation began by treating the choice of edits as a search over a whole dataset. AutoAugment searches over operations, probabilities and magnitudes with validation accuracy as the reward~\cite{cubuk2019auto}; because this search is expensive, RandAugment reduces it to two shared parameters~\cite{cubuk2020rand}, and TrivialAugment goes further, showing that sampling one edit and strength uniformly at random is a strong baseline that needs no tuning~\cite{muller2021trivial}. Detection raises an additional question, namely where in the image an edit should act. Zoph \emph{et al.} answer it by learning policies that include edits restricted to the content of each box~\cite{zoph2020det}, Scale-Aware AutoAug by conditioning image- and box-level policies on object scale~\cite{chen2021scale}, and InterAug by constructing a tuning-free detection policy~\cite{thopalli2023interaug}. Inspection work has followed the same pattern with recipes designed for damage imagery, such as bounding-box augmentation combined with samples from generative adversarial networks~\cite{aghayan2024} or content- and perspective-aware cut-and-paste of damage~\cite{rdd_cutpaste}. Although these methods differ in how the policy is found, they share one property: the edit for an image is drawn from a single distribution, conditioned at most on location or scale, which implicitly assumes that every class benefits from the same edits.

\noindent\textbf{Class-aware augmentation.} Evidence from recognition suggests that this assumption does not always hold. Kirichenko \emph{et al.} show that augmentation which raises average ImageNet accuracy can lower the accuracy of individual classes~\cite{kirichenko2023}, and Tsuchiya \emph{et al.} find that augmentation methods for road-damage classification differ in their effect on individual damage classes~\cite{tsuchiya2019}. Several methods have therefore made augmentation depend on the class. Yoo and Kang derive class-specific parameters from a classifier trained beforehand~\cite{yoo2023}, AdaAug learns class-adaptive policies in a separate search~\cite{cheung2022adaaug}, CUDA updates a severity for each class during training~\cite{ahn2023cuda}, and Pham \emph{et al.} set augmentation probabilities for each damage type from an exploration of the data~\cite{rdd_detectron2}. These methods share a limitation that becomes decisive in detection. The classification methods select the augmentation by the label of a sample, which has no answer when an image contains objects of several classes, while the detection recipe sidesteps the question by fixing its probabilities before training. PSIS, the closest detection counterpart, adapts to per-class detection accuracy during training, but it controls how many instance-switched samples each class receives rather than which edit each image receives~\cite{wang2019psis}. A method for multi-class detection must therefore choose edits from the whole set of classes in an image, and must do so as the needs of the detector change during training.

\noindent\textbf{Adapting augmentation during training.} Online methods make this second requirement possible by updating augmentation from the model as it trains, but they differ in the feedback they treat as evidence that an edit is useful. Recognition-based methods ask whether the current model still recognizes augmented examples. CTAugment adjusts per-operation magnitudes from predictions on augmented labeled data, pooled over classes~\cite{berthelot2020}; CUDA raises or lowers the severity of each class according to how well its augmented examples are recognized~\cite{ahn2023cuda}; and TeachAugment keeps learned transformations recognizable to a teacher while making them hard for the target model~\cite{suzuki2022teach}. Recognition, however, reflects what the model can already handle rather than what it gains from training on the edit, which is the quantity augmentation is meant to increase. Benefit-based methods measure this gain more directly. BDA rewards an augmentation by the training gain it yields over a paired unaugmented reference~\cite{lu2023bda}, DDAS uses a one-step gradient update to estimate how the augmentation distribution affects a held-out loss, including for detection~\cite{liu2021ddas}, MADAug learns an input-dependent policy through a bilevel validation objective~\cite{hou2023madaug}, and RangeAugment and AdaAugment adapt magnitudes from task-loss feedback~\cite{mehta2022rangeaugment,yang2025adaaugment}. Their common limitation is that the benefit is computed over all classes together, so a gain for frequent classes can outweigh a loss for a rare one, which is precisely the trade-off that a multi-class image creates. Estimating the gain separately for each class need not be costly. TracIn measures the influence of a training example on a test loss through the alignment of their gradients~\cite{pruthi2020tracin}, and CUA applies the same first-order argument to edited versus unedited images and to a separate reference loss for each class.

